% §1.2 Quantifying the Noise
%
% Cross-reference resolution status:
%
% ALREADY DEFINED (within chapter 1, by this section):
%   \label{sec:quantifying-noise}           -- this section
%   \label{tab:contingency}                 -- 2x2 contingency table
%
% RESOLVED:
%   \Cref{sec:three-faces}                  -- §1.1
%   \Cref{prop:complement-swap}             -- chapter 6 §6.2 (set-level complement)
%   \Cref{sec:bsc}                          -- chapter 2 §2.1
%   \Cref{ch:classification-measures}       -- chapter 7
%
% FORWARD REFS to later chapter 1 sections:
%   \Cref{sec:boolean-universe-first-look}  -- §1.3
%   \Cref{sec:sampling-rates}               -- §1.4
%
% Monograph register: fixes the contingency-table vocabulary used through
% Part I; the precision/recall cross-walk carries the prevalence point.
% Two remarks salvage former exercises (NOT swaps the rate pair; the
% constant noisy bit and the zero-information line).

\section{Quantifying the Noise}
\label{sec:quantifying-noise}

A noisy bit is fixed by two conditional error rates. We make ``rate'' precise
through the $2 \times 2$ contingency table, name the four rates it yields, and
relate them to the precision and recall of binary classification.

\subsection*{The contingency table}

Run a noisy bit $n$ times. Each trial has a known true input, $\top$ or $\bot$,
and an observed output, so the $n$ trials are ordered pairs (truth, output),
each falling into one of four joint states: a \emph{true positive} (truth
$\top$, output $\top$), a \emph{false negative} (truth $\top$, output $\bot$),
a \emph{false positive} (truth $\bot$, output $\top$), and a \emph{true
negative} (truth $\bot$, output $\bot$). The term ``false negative'' names the
output, not the truth: a $\bot$ reported when the truth was $\top$. Writing TP,
FN, FP, TN for the counts, so that $\mathrm{TP} + \mathrm{FN} + \mathrm{FP} +
\mathrm{TN} = n$, arrange them by truth (rows) and observed output (columns):

\begin{table}[ht]
\centering
\begin{tabular}{l|cc|c}
\hline
              & \textbf{output $\top$} & \textbf{output $\bot$} & total \\
\hline
\textbf{truth $\top$} & TP & FN & TP + FN \\
\textbf{truth $\bot$} & FP & TN & FP + TN \\
\hline
total         & TP + FP & FN + TN & $n$ \\
\hline
\end{tabular}
\caption{The $2 \times 2$ contingency table for $n$ observations of a noisy
bit. Rows condition on truth; the diagonal cells record agreement and the
off-diagonal cells disagreement.}
\label{tab:contingency}
\end{table}

\subsection*{Rates}

The two error rates are conditional, each normalized by its truth row:
\[
  \fprate \;=\; \frac{\mathrm{FP}}{\mathrm{TN} + \mathrm{FP}},
  \qquad
  \fnrate \;=\; \frac{\mathrm{FN}}{\mathrm{TP} + \mathrm{FN}}.
\]
Dividing by the row total rather than by $n$ is what makes $\fprate$ and
$\fnrate$ properties of the noisy bit itself rather than of the input mixture
observed. The complementary rates are the true-negative rate
$\tnrate = 1 - \fprate$ (the $\bot$-inputs reported correctly, also
\emph{specificity}) and the true-positive rate $\tprate = 1 - \fnrate$ (also
\emph{sensitivity}, and \emph{recall} in the vocabulary below). The two free
parameters of a noisy bit are $\fprate$ and $\fnrate$; the other two are
derived. The biased coin's $(\tau, \epsilon)$ of \Cref{sec:three-faces} is
$(\tprate, \fprate)$, and the binary symmetric channel of \Cref{sec:bsc} is the
constrained case $\fprate = \fnrate$, where one parameter suffices. These four
symbols keep these meanings throughout the book.

\begin{remark}[Negation reflects the rate pair]
\label{rem:not-swaps-rates}
Relabeling the output of a noisy bit, reporting $\bot$ wherever it reported
$\top$ and conversely, exchanges the two error directions: a noisy bit with
rates $(\fprate, \fnrate)$ becomes one with rates $(\fnrate, \fprate)$. Logical
negation thus acts on the parameter pair as the reflection
$(\fprate, \fnrate) \mapsto (\fnrate, \fprate)$ across the diagonal of the rate
square $[0,1]^2$. This is the atom-level form of the complement operation that
Part~II lifts to approximate sets (\Cref{prop:complement-swap}).
\end{remark}

\subsection*{Precision and recall}

A practitioner of binary classification reads \Cref{tab:contingency} through
precision and recall:
\[
  \mathrm{precision} \;=\; \frac{\mathrm{TP}}{\mathrm{TP} + \mathrm{FP}},
  \qquad
  \mathrm{recall} \;=\; \frac{\mathrm{TP}}{\mathrm{TP} + \mathrm{FN}}.
\]
Recall coincides with the true-positive rate, $\mathrm{recall} = \tprate = 1 -
\fnrate$. Precision does not reduce to the rates of the noisy bit: its
denominator is a column total, which depends on how many $\top$- and
$\bot$-inputs the sample contained. With $n_\top$ positives and $n_\bot$
negatives, $\mathrm{TP} \approx n_\top \tprate$ and $\mathrm{FP} \approx n_\bot
\fprate$, so
\[
  \mathrm{precision}
  \;\approx\;
  \frac{n_\top (1 - \fnrate)}{n_\top (1 - \fnrate) + n_\bot \fprate}.
\]
Precision is therefore a function of the \emph{prevalence} $n_\top / n_\bot$ as
well as the rates, whereas $(\fprate, \fnrate)$ is intrinsic to the noisy bit.
The pair $(\fprate, \fnrate)$ describes the noisy bit; the pair (precision,
recall) describes the noisy bit together with a workload. This is why the
framework takes $(\fprate, \fnrate)$ as the primary description and treats
precision and accuracy as derived; \Cref{ch:classification-measures} develops
the derived measures as random variables.

\subsection*{An instance}

A filter run on $n = 100$ messages, 40 spam ($\top$) and 60 ham ($\bot$),
produces the table

\begin{center}
\begin{tabular}{l|cc|c}
\hline
              & \textbf{flagged spam} & \textbf{flagged ham} & total \\
\hline
\textbf{true spam} & 35 & 5  & 40  \\
\textbf{true ham}  & 3  & 57 & 60  \\
\hline
total         & 38 & 62 & 100 \\
\hline
\end{tabular}
\end{center}

\noindent
with rates $\fprate = 3/60 = 0.05$ and $\fnrate = 5/40 = 0.125$, hence
$\tnrate = 0.95$ and $\tprate = 0.875$. Recall is $35/40 = 0.875 = \tprate$,
while precision is $35/38 \approx 0.921$, exceeding $\tprate$ because the column
total (38) falls short of the truth row total (40). At a different prevalence
the same $(\fprate, \fnrate)$ would yield a different precision.

\begin{remark}[The constant noisy bit and the zero-information line]
\label{rem:zero-information-line}
A noisy bit that ignores its input, reporting $\top$ with a fixed probability
$p$ regardless of the truth, has $\fprate = p$ and $\fnrate = 1 - p$, so
$\fprate + \fnrate = 1$. On this anti-diagonal of the rate square the output is
statistically independent of the truth: the report carries no information about
the input. A constant rule is thus simultaneously a perfect classifier for one
truth class and the worst possible one for the other (at $p = 1$, $\fnrate = 0$
but $\fprate = 1$). The locus $\fprate + \fnrate = 1$ is the zero-information
case, and it is distinct from the symmetric case $\fprate = \fnrate$, which
remains informative for $\fprate \neq 1/2$.
\end{remark}
